\section{Conclusiones}

El análisis muestra que un modelo E--R no se limita a enumerar entidades y
relaciones: debe conservar la dependencia semántica de los atributos, hacer
explícitas las restricciones de identidad, cardinalidad y participación, y
documentar toda decisión añadida ante información incompleta. Las relaciones
ternarias, la agregación, las especializaciones, las categorías y las
relaciones recursivas resuelven necesidades diferentes y no deben sustituirse
entre sí sin comprobar qué información se perdería.

Los modelos obtenidos satisfacen los requisitos de cada mini--mundo y mantienen
coherencia entre los diagramas y su explicación. La entrega conjunta de las
fuentes editables, sus exportaciones legibles y el reporte permite verificar
tanto la corrección conceptual como las decisiones asumidas durante el diseño.
